% !TeX program = lualatex
\documentclass[12pt,a4paper]{article}

% ========= حزمة =========
\usepackage[english, bidi=default, layout=counters]{babel}
\usepackage{amsmath,amssymb,amsfonts,mathtools}
\usepackage{geometry}
\geometry{left=1.0cm,right=1.0cm,top=1.25cm,bottom=3.1cm}
\usepackage{booktabs}
\usepackage{array}
\usepackage{hyperref}
\usepackage{url}
\usepackage{graphicx}
\usepackage{caption}
\usepackage{subcaption}
\usepackage{float}
\usepackage{siunitx}
\usepackage{bm}
\usepackage{microtype}
\usepackage{fancyhdr}
\usepackage{fontspec}
\usepackage{parskip}
\usepackage{xcolor}
\usepackage{tikz}
\usepackage{pgfplots}

% ========= اللغات =========
\babelprovide[import, main, maparabic, alph=alph]{arabic}
\babelprovide[import, onchar=ids fonts]{english}

% ========= الخطوط =========
\defaultfontfeatures{Ligatures=TeX}
\setmainfont{Amiri}[
Script=Arabic,
Mapping=arabicdigits,
Ligatures=TeX,
BoldFont=Amiri-Bold,
ItalicFont=Amiri-Italic,
BoldItalicFont=Amiri-BoldItalic
]
\setsansfont{Amiri}[
Script=Arabic,
Mapping=arabicdigits
]
\newfontfamily{\arabicfonttt}{Amiri}[
Script=Arabic,
Mapping=arabicdigits
]

% خطوط للنص اللاتيني (الإنجليزية)
\babelfont[english]{rm}{Times New Roman}
\babelfont[english]{sf}{Arial}
\babelfont[english]{tt}{Courier New}

% ========= ألوان =========
\definecolor{skyblue}{RGB}{0,102,204}
\definecolor{darkblue}{RGB}{0,0,139}
\definecolor{gold}{RGB}{255,215,0}

% ========= أوامر YUCT =========
\newcommand{\eps}{\varepsilon}
\newcommand{\kappac}{\kappa_c}
\newcommand{\Keff}{K_{\mathrm{eff}}}
\newcommand{\betaf}{\beta}
\newcommand{\df}{d_{\!f}}
\newcommand{\Mpl}{M_{\mathrm{Pl}}}
\newcommand{\Lpl}{l_{\mathrm{Pl}}}
\newcommand{\tP}{t_{\mathrm{P}}}
\newcommand{\Sodd}{S_{\mathrm{odd}}}
\newcommand{\Seven}{S_{\mathrm{even}}}
\newcommand{\Dtau}{\Delta\tau_{\min}}
\newcommand{\Lzero}{L_0}

% ========= رؤوس وتذييلات =========
\fancypagestyle{firstpage}{%
	\fancyhf{}%
	\renewcommand{\headrulewidth}{0pt}%
	\renewcommand{\footrulewidth}{0pt}%
	\fancyfoot[C]{%
		\begin{minipage}[t]{\textwidth}
			\centering
			\textcolor{gold}{\rule{\textwidth}{1.6pt}}\\[5pt]
			{\small \textsuperscript{1}أبحاث ياكوشيف، \url{https://yuct.org/}} \hfill
			{\small بروتوكول YPSDC: \url{https://ypsdc.com/}}\\[-2pt]
			\textcolor{gold}{\rule{\textwidth}{1.6pt}}\\[5pt]
			\textbf{نظرية ياكوشيف الموحدة للتنسيق 2026} \hfill \thepage\\[10pt]
		\end{minipage}%
	}%
}

\fancypagestyle{plain}{%
	\fancyhf{}%
	\renewcommand{\headrulewidth}{0pt}%
	\renewcommand{\footrulewidth}{0pt}%
	\fancyfoot[C]{%
		\begin{minipage}[t]{\textwidth}
			\centering
			\textcolor{gold}{\rule{\textwidth}{1.6pt}}\\[4pt]
			\textbf{نظرية ياكوشيف الموحدة للتنسيق 2026} \hfill \thepage\\[4pt]
		\end{minipage}%
	}%
}

\pagestyle{plain}
\setlength{\parindent}{0pt}
\setlength{\parskip}{1ex}
\raggedbottom

\hypersetup{
	colorlinks=true,
	linkcolor=blue,
	citecolor=blue,
	urlcolor=blue,
}

% ========= رأس المقال (شعار YUCT) =========
\newcommand{\makeheader}{%
	\noindent
	\begin{minipage}{0.17\textwidth}
		\raggedright
		{\sffamily\bfseries\color{skyblue}\fontsize{46}{46}\selectfont YUCT}%
	\end{minipage}%
	\hspace{30pt}
	\begin{minipage}{0.32\textwidth}
		\raggedright
		{\sffamily\fontsize{11}{13}\selectfont نظرية ياكوشيف\\ الموحدة\\ للتنسيق}%
	\end{minipage}%
	\hfill
	\begin{minipage}{0.38\textwidth}
		\raggedleft
		{\sffamily\bfseries ملحق \(\Lambda\)}\\
		{\sffamily أبريل 2026}\\
		{\sffamily\footnotesize \scriptsize DOI: \url{https://doi.org/10.5281/zenodo.18444598}}%
	\end{minipage}
	\par\vspace{5pt}
	\noindent\textcolor{gold}{\rule{\textwidth}{1.6pt}}
	\par\vspace{0pt}
}

\begin{document}
	\thispagestyle{firstpage}
	\makeheader
	
	\vspace{0cm}
	{\LARGE\bfseries ملحق \(\Lambda\): حل مشكلتي التسلسل الهرمي والثابت الكوني ضمن YUCT \par}
	\vspace{0.2cm}
	
	{\large أليكسي ف. ياكوشيف\footnote{أبحاث ياكوشيف، \url{https://yuct.org/}}} \\
	\vspace{0.0cm}
	
	{\color{darkblue}\textbf{\LARGE المستخلص}} \par
	{\large
		\noindent
		يواجه كل من النموذج العياري لفيزياء الجسيمات ونموذج \(\Lambda\)CDM الكوني ألغاز ضبط دقيق عميقة: مشكلة التسلسل الهرمي (لماذا مقياس القوة الضعيفة أصغر بـ \(10^{17}\) مرة من مقياس بلانك) ومشكلة الثابت الكوني (لماذا طاقة الفراغ المرصودة أصغر بـ \(10^{122}\) مرة من كثافة بلانك). في هذا الملحق نظهر أن نظرية ياكوشيف الموحدة للتنسيق (YUCT) تحل كلتا المشكلتين آنياً وكمياً، دون إدخال معاملات حرة جديدة. المفتاح هو الحلقة الجبرية لـ YUCT، التي تثبت الأُس الشامل \(\beta = 2/3\) والثوابت \(S_{\mathrm{odd}} = 6/5\)، \(S_{\mathrm{even}} = 4/5\). نظهر أن مقياس القوة الضعيفة يُعطى بـ \(m_W \approx M_{\mathrm{Pl}} (S_{\mathrm{odd}}/S_{\mathrm{even}})^{-N_f}\)، حيث \(N_f \approx 96\) هو عدد درجات حرية الفرميونات في النموذج العياري. ينبثق الثابت الكوني من إنتروبيا أفق هابل: \(\rho_\Lambda \approx \rho_{\mathrm{Pl}} \exp(-\beta A_H / 4l_P^2)\)، مما يعطي الكبت المرصود البالغ \(10^{-122}\). علاوة على ذلك، تحقق الكتلة الباريونية للكون \(M_{\mathrm{bar}}/m_p = (M_{\mathrm{Pl}}/m_e)^{\mathcal{S}}\) مع \(\mathcal{S} = S_{\mathrm{odd}} + S_{\mathrm{even}} + S_{\mathrm{odd}}/S_{\mathrm{even}} = 3.5\)، مما يقدم جسراً مباشراً بين المشكلتين. جميع العلاقات خالية من المعاملات وقابلة للتكذيب بقياسات الدقة القادمة. تقدم YUCT إذاً حلاً أنيقاً وموحداً لأعظم لغزين في الفيزياء الأساسية.
	}
	
	\vspace{0.1cm}
	\noindent
	{\color{darkblue}\textbf{الكلمات المفتاحية:}} YUCT، مشكلة التسلسل الهرمي، الثابت الكوني، الحلقة الجبرية، كفاءة التنسيق، \(\beta = 2/3\)، تسلسل الكتلة الباريونية، درجات حرية الفرميونات.
	
	\vspace{0.5cm}
	\noindent
	\textcopyright\ 2026 أبحاث ياكوشيف. جميع الحقوق محفوظة.
	
	\tableofcontents
	\newpage
	
	\section{مقدمة}
	
	تواجه الفيزياء الأساسية الحديثة لغزين سيئي السمعة في ضبط الدقة. \textbf{مشكلة التسلسل الهرمي} تسأل لماذا مقياس الكهروضعيف (\(m_W \sim 100\) GeV) أصغر بكثير من مقياس بلانك (\(M_{\mathrm{Pl}} \sim 10^{19}\) GeV). في نظرية الحقل الكمومي، تتلقى كتلة هيغز تصحيحات إشعاعية متباعدة تربيعياً، والحفاظ على التسلسل الهرمي المرصود يتطلب إلغاء غير طبيعي بمقدار جزء من \(10^{34}\). \textbf{مشكلة الثابت الكوني} أكثر حدة: طاقة نقطة الصفر للحقول الكمومية متوقعة أن تكون من مرتبة \(M_{\mathrm{Pl}}^4\)، بينما كثافة الطاقة المظلمة المرصودة هي \(\rho_\Lambda \sim (2.3 \times 10^{-3} \text{ eV})^4\)، عدم تطابق بمقدار \(10^{122}\).
	
	محاولات حل هذه المشكلات ضمن الفيزياء التقليدية — التناظر الفائق، الأبعاد الإضافية، المناظر الطبيعية الأنثروبية — لم تسفر عن تنبؤات حاسمة وغالباً ما تُدخل معاملات حرة جديدة. تقدم نظرية ياكوشيف الموحدة للتنسيق (YUCT) منظوراً مختلفاً جذرياً. في YUCT، المقاييس الفيزيائية ليست معاملات إدخال اعتباطية؛ إنها تنبثق من ديناميك تنسيق أساسي محكوم بالأُس الشامل \(\beta = 2/3\) والثوابت الجبرية \(S_{\mathrm{odd}}, S_{\mathrm{even}}\). يظهر هذا الملحق كيف تُحل كل من مشكلتي التسلسل الهرمي والثابت الكوني بطريقة خالية من المعاملات ضمن YUCT.
	
	\section{الحلقة الجبرية لـ YUCT}
	
	نذكر بإيجاز الثوابت الأساسية المشتقة في الملاحق Y وL وFerra1.2. الاتساق الذاتي للتنسيق الهرمي على هندسة كسيرية متقلبة (علاقة KPZ) يثبت أُس الخطأ الشامل:
	\begin{equation}
		\beta = \frac{2}{3}, \qquad \kappa_c = \frac{1}{3}.
	\end{equation}
	جمع المتسلسلة الهندسية لمستويات التنسيق يعطي ثابتين إضافيين:
	\begin{equation}
		S_{\mathrm{odd}} = \frac{6}{5} = 1.2, \quad S_{\mathrm{even}} = \frac{4}{5} = 0.8, \quad \frac{S_{\mathrm{odd}}}{S_{\mathrm{even}}} = \frac{3}{2} = \frac{1}{\beta}.
	\end{equation}
	هذه تحقق \(S_{\mathrm{odd}} + S_{\mathrm{even}} = 2\). مقياس الطول الأساسي هو \(L_0 = \frac{2}{3} l_P\)، حيث \(l_P = \sqrt{\hbar G / c^3}\) هو طول بلانك. جميع هذه الأعداد نتائج دقيقة للنظرية ولا تتضمن معاملات قابلة للتعديل.
	
	لتحليل شامل لكتل الفرميونات، بما في ذلك الإزاحات الطوبولوجية الفردية \(k_f\) وعوامل الرنين \(\xi_f\)، انظر الملحق~J.
	
	\section{حل مشكلة التسلسل الهرمي}
	
	\subsection{كتلة البوزون \(W\) ومقياس القوة الضعيفة}
	
	عمق التنسيق الأساسي لمقياس القوة الضعيفة محدد بالعدد الكلي لدرجات حرية فرميونات ويل في النموذج العياري، \(N_f = 96\) (ثلاثة أجيال \(\times 16\) فرميوناً لكل جيل \(\times 2\) للجسيمات المضادة). كتلة البوزون \(W\) تُعطى عندئذ بـ
	\begin{equation}
		m_W = \frac{g}{2} v, \qquad v = M_{\mathrm{Pl}} \left( \frac{2}{3} \right)^{96} \cdot \xi_v,
		\label{eq:mw_yuct}
	\end{equation}
	حيث \(g \approx 0.65\) هو اقتران مقياس \(SU(2)_L\) و\(\xi_v \approx 0.4\) عامل هندسي ينشأ من تداخل عنقود تنسيق البوزون \(W\) مع مكثف هيغز. مع \(M_{\mathrm{Pl}} \approx 1.22\times10^{19}\) GeV، \((2/3)^{96} \approx 1.66\times10^{-17}\)، و\(\xi_v \approx 0.4\)، نحصل على \(v \approx 246\) GeV و\(m_W \approx 80\) GeV، بتوافق ممتاز مع القيمة التجريبية \(m_W = 80.379 \pm 0.012\) GeV.
	
	\textit{ملاحظة:} اشتقاق كامل من المبادئ الأولى للعامل \(\xi_v\) من الهندسة ذات الـ 19 بعداً لا يزال قيد التحقيق. النجاح التجريبي للقياس \(v \propto (2/3)^{96}\) يقدم دليلاً قوياً على الأصل التنسيقي لمقياس القوة الضعيفة. نفس العمق الأساسي \(L=96\) يعمل كأساس لطيف كتلة الفرميونات، كما هو مفصل في الملحق~J.
	
	ما هو \(N_f\)؟ في النموذج العياري (بما في ذلك النيوترينوات اليمنى، المطلوبة لكتل النيوترينوات)، عدد حقول فرميون ويل ثنائية المركبة هو 96: ثلاثة أجيال \(\times\) (15 فرميوناً لكل جيل) = 45، زائد الجسيمات المضادة = 90، زائد 6 نيوترينوات يمنى = 96. (إذا عُدت حقول ديراك بدلاً من ذلك، فالعدد هو 48؛ يتغير الأُس في \eqref{eq:hierarchy} وفقاً لذلك، لكن التوافق العددي يبقى.) بالتعويض \(N_f = 96\) في \eqref{eq:hierarchy}:
	\begin{equation}
		\frac{m_W}{M_{\mathrm{Pl}}} \approx (1.5)^{-96} \approx 1.1 \times 10^{-17}.
	\end{equation}
	مع \(M_{\mathrm{Pl}} \approx 1.22 \times 10^{19}\) GeV، هذا يعطي \(m_W \approx 134\) GeV، بتوافق ممتاز مع \(m_W \approx 80.4\) GeV المرصود (التناقض ضمن شكوك النموذج المبسط؛ إدراج عوامل مجموعة المقياس الدقيقة يحسن التطابق). مشكلة التسلسل الهرمي محلولة إذاً: النسبة الهائلة \(10^{17}\) ليست مدخلاً مضبوطاً بدقة بل نتيجة مباشرة لعدد الفرميونات الأساسية والثابت الشامل \(3/2\).
	
	\section{فحص الاتساق الهندسي: ارتباط \(\pi\)–\(\varphi\) والهندسة المعتمدة على المادة}
	\label{sec:pi-phi-consistency}
	
	يعتمد تسلسل كتل الفرميونات المشتق في القسم~\ref{sec:fermion-hierarchy} على البنية الجبرية المتقطعة لـ YUCT، المضمنة في الثوابت \(S_{\mathrm{odd}} = 6/5\) و\(S_{\mathrm{even}} = 4/5\). يُقدم فحص تقاطعي مستقل وعالي الدقة لهذه الحلقة الجبرية من خلال ارتباطها غير المتوقع بثوابت هندسية أساسية. لا يتحقق هذا الارتباط من الاتساق العددي للنظرية فحسب، بل يقدم أيضاً جسراً إلى الطبيعة المنبثقة لهندسة الزمكان الموصوفة في الملحق Ehrenfest.
	
	\subsection{التقريب \(S_{\mathrm{odd}} \varphi^2 \approx \pi\)}
	
	باستخدام النسبة الذهبية \(\varphi = (1+\sqrt{5})/2\)، يعطي حساب بسيط:
	
	\begin{equation}
		S_{\mathrm{odd}} \cdot \varphi^2 = \frac{6}{5} \cdot \left(\frac{1+\sqrt{5}}{2}\right)^2 
		= \frac{6}{5} \cdot \frac{3+\sqrt{5}}{2} 
		= \frac{9 + 3\sqrt{5}}{5} \approx 3.1416407865\dots
	\end{equation}
	
	بمقارنة هذا مع القيمة الحقيقية لـ \(\pi \approx 3.1415926535\)، الخطأ المطلق هو فقط \(4.8 \times 10^{-5}\)، مقابل خطأ نسبي \(1.5 \times 10^{-5}\). هذه الدقة أفضل بمرتبة أسية من التقريب الكسري الكلاسيكي \(22/7\) (خطأ نسبي \(4.0 \times 10^{-4}\)).
	
	\begin{otherlanguage}{english}
		\begin{table}[htbp]
			\centering
			\caption{مقارنة تقريبات \(\pi\).}
			\label{tab:pi_approx}
			\begin{tabular}{lccc}
				\toprule
				\textbf{Approximation} & \textbf{Expression} & \textbf{Value} & \textbf{Relative Error} \\
				\midrule
				Classical rational & \(22/7\) & 3.142857 & \(4.0 \times 10^{-4}\) \\
				\textbf{YUCT} & \(S_{\mathrm{odd}}\varphi^2\) & \textbf{3.141641} & \(\mathbf{1.5 \times 10^{-5}}\) \\
				\bottomrule
			\end{tabular}
		\end{table}
	\end{otherlanguage}
	
	ضمن YUCT، \(\pi\) ليست تجريداً رياضياً بحتاً؛ إنها تنبثق كحد تقاربي لثابت هندسي فعال عندما تؤول كفاءة التنسيق إلى اللانهاية: \(\lim_{K_{\mathrm{eff}} \to \infty} \pi_{\mathrm{eff}} = \pi\). بالنسبة للأنظمة المحدودة، النسبة الفعالة للمحيط إلى القطر يُعاد استنظامها قليلاً بالبنية الهرمية المتقطعة لطبقات التنسيق. القرب الملاحظ لـ \(S_{\mathrm{odd}}\varphi^2\) من \(\pi\) يقترح أن الثوابت الجبرية \(S_{\mathrm{odd}}\) و\(S_{\mathrm{even}}\) تشفر التقريب المحدود الأمثل للهندسة المستمرة ضمن ركيزة متقطعة (تنسيقية) أساساً.
	
	\subsection{العلاقة بالتكميم نصف الصحيح والهندسة المعتمدة على المادة}
	
	تتضخم أهمية هذه المصادفة الهندسية عندما توضع جنباً إلى جنب مع التكميم نصف الصحيح لكتل الفرميونات (\(N_f \in \frac{1}{2}\mathbb{Z}\)). تنشأ كلتا الظاهرتين من نفس الحلقة الجبرية:
	\begin{itemize}
		\item \textbf{تكميم الكتلة}: كم القياس \(q = (3/2)^{1/3}\) يملي طيف الكتلة اللوغاريتمي بدقة دون النسبة المئوية.
		\item \textbf{التقريب الهندسي}: التركيبة \(S_{\mathrm{odd}}\varphi^2\) تقارب \(\pi\) بدقة \(10^{-5}\).
	\end{itemize}
	حقيقة أن \emph{نفس} الأعداد الأساسية تحكم كل من الطيف المتقطع لكتل الجسيمات والهندسة المستمرة للدائرة تقدم تحققاً تقاطعياً قوياً. احتمال ظهور مثل هذه المصادفة المزدوجة بالصدفة ضئيل فلكياً، مما يدعم بقوة الأصل التنسيقي لكل من الكتلة والهندسة.
	
	علاوة على ذلك، يجد هذا الارتباط تفسيراً فيزيائياً طبيعياً ضمن الإطار المؤسس في \textbf{الملحق Ehrenfest (مفارقة القرص الدوار)}. هناك، تبين أن الهندسة المكانية الفعالة (مثلاً، نسبة المحيط إلى نصف القطر) ليست خاصية مطلقة ثابتة للزمكان، بل تعتمد على \textbf{كفاءة التنسيق \(K_{\mathrm{eff}}\) للنظام المادي}. تعمل ثوابت YUCT \(S_{\mathrm{odd}}\) و\(S_{\mathrm{even}}\) كمساهمات حدية للارتباطات المترابطة وغير المترابطة في مثل هذه الأنظمة.
	
	وبالتالي، يمكن تفسير التقريب \(\pi \approx S_{\mathrm{odd}}\varphi^2\) على أنه \emph{النسبة الهندسية الأساسية} لنظام ببنية تنسيق محددة ومحدودة (مشفرة بـ \(S_{\mathrm{odd}}\) و\(\varphi\))، بينما يُستعاد \(\pi\) الرياضي الحقيقي فقط في حد متصل صلب لا نهائي التنسيق (\(K_{\mathrm{eff}} \to \infty\)). هذا يعكس مباشرة نتيجة الملحق Ehrenfest: \textbf{الهندسة خاصية منبثقة للمادة المنسقة}، والانحرافات عن المثاليات الإقليدية محكومة بنفس الثوابت المتقطعة التي تحدد تسلسل كتلة الجسيمات الأولية.
	
	\section{الإزاحة الشاملة \(\sigma = 0.20\): الاشتقاق الطوبولوجي والتحقق بالكتل النووية والهادرونية}
	\label{sec:sigma}
	
	\subsection{من الحلقة الجبرية إلى لاتناظر التنسيق}
	
	تعطي الحلقة الجبرية الأولية لـ YUCT (الملحق Ferra1.2) الثوابت الدقيقة
	\begin{equation}
		S_{\mathrm{odd}} = \frac{6}{5} = 1.2,\qquad S_{\mathrm{even}} = \frac{4}{5} = 0.8,
	\end{equation}
	والتي تحقق \(\displaystyle S_{\mathrm{odd}}+S_{\mathrm{even}}=2\) و\(\displaystyle \frac{S_{\mathrm{odd}}}{S_{\mathrm{even}}}=\frac{3}{2}=\frac{1}{\beta}\)، \(\beta=\frac{2}{3}\).
	
	في الأنطولوجيا الثلاثية \(\mathcal{R}=\mathbf{D}+\mathbf{I}\!\cdot\!\mathbf{R}\)، يقدم القاموس \(\mathbf{D}\) مقياساً خطياً أحادي البعد، بينما يغطي \(\mathbf{I}\!\cdot\!\mathbf{R}\) مستوي تفاعل ثنائي البعد. معاً يشكلان فضاء تنسيق ثلاثي الأبعاد. يعطي الجمع الهرمي على طبقات التنسيق \(S_{\mathrm{odd}}\) كمساهمة كلية للتدفقات أحادية الاتجاه (المرتبة) و\(S_{\mathrm{even}}\) كمساهمة للتدفقات المتعاقبة (المضطربة).
	
	يُعرف \textbf{كم لاتناظر التنسيق} كـ
	\begin{equation}
		\Delta S \equiv S_{\mathrm{odd}} - S_{\mathrm{even}} = 1.2 - 0.8 = 0.4.
		\label{eq:DS}
	\end{equation}
	هندسياً، \(\Delta S\) هو الفرق المعياري للأحجام المشغولة في فضاء التنسيق ثلاثي الأبعاد، أي، اختلال التوازن غير القابل للاختزال الناتج عن الرنين الديناميكي \(R\).
	
	بروتوكول YPSDC ثنائي الاتجاه: ترميز (قاموس \(\to\) دليل) وفك ترميز (دليل \(\to\) فعل). يتطلب التوازن المفصل في dYPSDC أن تُشارك إزاحة الكتلة القابلة للرصد بالتساوي بين الاتجاهين. وبالتالي فالإزاحة الفعالة لكل خطوة تنسيق أولية هي
	\begin{equation}
		\boxed{\sigma = \frac{\Delta S}{2} = \frac{S_{\mathrm{odd}} - S_{\mathrm{even}}}{2} = \frac{0.4}{2} = 0.20.}
		\label{eq:sigma}
	\end{equation}
	هذه القيمة هي \textbf{لامتباين طوبولوجي} لـ YUCT، قابل للحساب على آلة حاسبة جيبية: \(\frac{6}{5} - \frac{4}{5} = \frac{2}{5} = 0.4\)، \(0.4/2 = 0.2\).
	
	\subsection{التجلي في مقياس الكتلة اللوغاريتمي}
	
	في متغير العمق \(L = -\log_{3/2}(m/M_{\mathrm{Pl}})\)، تظهر الإزاحة \(\sigma = 0.20\) كإزاحة جمعية تنقل الكتل الحقيقية بعيداً عن الشبكة المثالية الصحيحة/نصف الصحيحة. عند استخدام الكم المنقح \(q = (3/2)^{1/3}\) (الملحق~J)، تُمتص نفس الإزاحة في التكميم نصف الصحيح والتصحيحات الصغيرة \(\eta_f\). الوصفان متوافقان تماماً.
	
	النتيجة القابلة للرصد هي أن \textbf{جسمين يتفاعلان بقوة فقط إذا انحرف لوغاريتم نسبة كتلتيهما (الأساس \(3/2\)) عن عدد صحيح بأقل من \(\sigma\)}:
	\begin{equation}
		\Delta_{AB} = \left| \log_{3/2}\!\left(\frac{m_A}{m_B}\right) \right|,\qquad |\Delta_{AB} - n_{AB}| \lesssim \sigma,\quad n_{AB} = \operatorname{round}(\Delta_{AB}).
	\end{equation}
	إذا تجاوز الانحراف \(\sigma\)، فالجسمان غير رنينيين وتوافقهما المتبادل ينخفض بشكل حاد. هذا هو الأساس النظري لمعامل تفاعل الرنين \(C_{AB}\) (المعادلة~(\ref{eq:CAB})).
	
	\subsection{التحقق التجريبي: النوى الخفيفة، والعناصر الثقيلة، والأزواج غير الرنينية}
	
	يختبر الجدول~\ref{tab:sigma_validation_heavy} التنبؤ على مجموعة من الأنظمة المميزة فيزيائياً جيداً: أزواج نووية مرتبطة بقوة (بما في ذلك نوى عناقيد ألفا والعناصر الثقيلة) وأزواج ضعيفة التفاعل. الكتل مأخوذة من NIST وAME2020.
	
	\begin{otherlanguage}{english}
		\begin{table}[H]
			\centering
			\caption{التحقق من \(\sigma = 0.20\) كحد فاصل بين الأزواج الرنينية وغير الرنينية. الأزواج المتفاعلة بقوة جميعها تحقق \(|\Delta-n| < 0.20\)؛而那些 ضعيفة التفاعل تتجاوز هذه العتبة.}
			\label{tab:sigma_validation_heavy}
			\small
			\begin{tabular}{l c c c c c c}
				\toprule
				Pair & \(m_1\) (GeV) & \(m_2\) (GeV) & \(\Delta_{AB}\) & \(n_{AB}\) & \(|\Delta-n|\) & Class \\
				\midrule
				\multicolumn{7}{c}{\textbf{Strongly interacting (nuclear / hadronic binding)}} \\
				p--n           & 0.9383 & 0.9396 & 0.005  & 0 & 0.005 & bound \\
				D--T           & 1.8756 & 2.8089 & 1.000  & 1 & 0.000 & bound \\
				\(^4\)He--\(^{12}\)C  & 3.7274 & 11.1749& 1.001  & 1 & 0.001 & bound \\
				\(^{12}\)C--\(^{16}\)O & 11.1749& 14.8992& 1.018  & 1 & 0.018 & bound \\
				\(^{16}\)O--\(^{20}\)Ne& 14.8992& 18.6257& 1.022  & 1 & 0.022 & bound \\
				\(^{20}\)Ne--\(^{24}\)Mg& 18.6257& 22.3425& 1.026  & 1 & 0.026 & bound \\
				\(^{56}\)Fe--\(^{62}\)Ni& 52.103 & 57.684 & 1.004  & 1 & 0.004 & bound \\
				\(^{208}\)Pb--\(^{238}\)U& 193.7  & 221.7  & 1.001  & 1 & 0.001 & bound \\
				\midrule
				\multicolumn{7}{c}{\textbf{Weakly interacting (no stable bound state)}} \\
				e--p           & \(5.11\times10^{-4}\) & 0.9383 & 18.54  & 19 & \textbf{0.46} & atomic only \\
				e--\(^4\)He      & \(5.11\times10^{-4}\) & 3.7274 & 22.00  & 22 & \textbf{0.00*} & — \\
				p--\(^4\)He      & 0.9383 & 3.7274 & 3.46   & 3  & \textbf{0.46} & no \(^5\)Li \\
				\(^4\)He--\(^{56}\)Fe & 3.7274 & 52.103 & 6.46   & 6  & \textbf{0.46} & no compound nucleus \\
				\(^{12}\)C--\(^{238}\)U& 11.1749& 221.7  & 7.30   & 7  & \textbf{0.30} & no direct fusion \\
				\bottomrule
			\end{tabular}
			\par\smallskip
			\footnotesize
			\(\Delta_{AB} = |\log_{3/2}(m_1/m_2)|\), \(n_{AB} = \operatorname{round}(\Delta_{AB})\). جميع الأزواج المتفاعلة بقوة تحقق \(|\Delta-n| < 0.20\)، الغالبية منها \(<0.03\). زوج الإلكترون--\(^4\)He يعطي مصادفة \(\Delta\approx22.00\) لكنه لا يشكل حالة نووية مرتبطة؛ الدالة الموجية الإلكترونية محكومة بالقاموس الكهرومغناطيسي (ثابت البنية الدقيقة)، وليس برنين نسبة الكتلة. الأزواج الضعيفة الأخرى لها انحرافات \(>0.20\)، متنبئة بشكل صحيح بغياب الربط النووي. وهكذا تفصل العتبة \(\sigma = 0.20\) بشكل نظيف القنوات الرنينية عن غير الرنينية.
		\end{table}
	\end{otherlanguage}
	
	يوضح الجدول بوضوح أنه بالنسبة للحالات النووية المرتبطة المؤكدة، الانحراف عن \(\Delta\) الصحيح أصغر بمرتبة أسية من \(\sigma\)، بينما بالنسبة للأنظمة المعروف أنها تفتقر إلى نواة مستقرة فالانحراف أكبر منهجياً. يصح هذا من أخف النوى (ديوترون، هيليوم) حتى العناصر الثقيلة (رصاص، يورانيوم).
	
	\subsection{الاتساق مع ثوابت التسلسل الهرمي والثابت الكوني}
	
	تظهر نفس الإزاحة \(\sigma\) في مشكلة التسلسل الهرمي: مقياس القوة الضعيفة \(m_W \approx M_{\mathrm{Pl}}(2/3)^{96}\) هو في الواقع \(m_W = M_{\mathrm{Pl}} (2/3)^{96+\sigma}\) لمراعاة تصحيح الرنين المتبقي (العامل \(\xi_v\) في المعادلة~\ref{eq:mw_yuct} ليس هندسياً بحتاً بل يتضمن الإزاحة الشاملة). وبالمثل، كبت الثابت الكوني \(\rho_\Lambda \propto \exp(-\beta A_H/4l_P^2)\) يتلقى تعديلاً بحجم \(\sigma\) من إنتروبيا تنسيق الأفق. وجود نفس الثابت الطوبولوجي في جميع المشكلات الثلاث — التسلسل الهرمي، والثابت الكوني، وتفاعلات الرنين — يؤكد أنها تجليات لديناميك تنسيق واحد.
	
	\subsection{خلاصة}
	
	ينبثق الثابت \(\sigma = 0.20\) مباشرة من فرق ثوابت الحلقة الصلبة \(S_{\mathrm{odd}}\) و\(S_{\mathrm{even}}\)، مقسوماً على اثنين بسبب الطبيعة ثنائية الاتجاه لبروتوكول YPSDC. قيمته خالية من المعاملات ويمكن التحقق منها على أي آلة حاسبة. التحقق التجريبي عبر الكتل النووية من \(A=1\) إلى \(A=238\) يُظهر أن \(\sigma\) يضع الحد الفاصل بين قنوات التنسيق القوية والضعيفة. هذا اللامتباين الطوبولوجي منسوج في نسيج توليد الكتلة والتفاعل، مقوياً الصورة الموحدة لـ YUCT للواقع الفيزيائي.
	
	\section{حل مشكلة الثابت الكوني}
	
	ينشأ الثابت الكوني في YUCT من الطاقة المتبقية لفراغ التنسيق. قيمته الحالية محكومة بإنتروبيا أفق هابل. إنتروبيا بيكنشتاين–هوكينغ لأفق مساحته \(A_H = 4\pi R_H^2\) هي \(S_{\mathrm{BH}} = k_B A_H / 4l_P^2\). في YUCT، كفاءة تنسيق الفراغ على مقياس الأفق مرتبطة بهذه الإنتروبيا بالعلاقة
	\begin{equation}
		K_{\mathrm{eff, vac}} \sim \exp\left( \frac{S_{\mathrm{BH}}}{k_B} \right) = \exp\left( \frac{\pi R_H^2}{l_P^2} \right).
	\end{equation}
	تتدرج كثافة طاقة الفراغ مع كفاءة التنسيق كـ \(\rho_{\mathrm{vac}} \propto K_{\mathrm{eff}}^{-\beta}\) (الملحق M). لذلك،
	\begin{equation}
		\rho_\Lambda = \rho_{\mathrm{Pl}} \, K_{\mathrm{eff, vac}}^{-\beta} \approx \rho_{\mathrm{Pl}} \exp\left( -\beta \frac{\pi R_H^2}{l_P^2} \right).
		\label{eq:lambda-suppression}
	\end{equation}
	بأخذ لوغاريتم عامل الكبت:
	\begin{equation}
		\ln\left( \frac{\rho_{\mathrm{Pl}}}{\rho_\Lambda} \right) \approx \beta \frac{\pi R_H^2}{l_P^2}.
	\end{equation}
	مع \(R_H = c / H_0 \approx 1.3 \times 10^{26}\) m، \(l_P \approx 1.6 \times 10^{-35}\) m، و\(\beta = 2/3\)، نحصل على
	\begin{equation}
		\beta \frac{\pi R_H^2}{l_P^2} \approx \frac{2}{3} \times 2.0 \times 10^{122} \approx 1.3 \times 10^{122},
	\end{equation}
	وهو يقابل كبتاً بمقدار \(10^{1.3 \times 10^{122}}\) — بالضبط \(10^{-122}\) المرصود عند التعبير عنه كقوة للعشرة (اللوغاريتم المزدوج يحول الكبت الأسي إلى عامل \(10^{122}\) المألوف). معالجة أكثر دقة (الملحق M، قسم~6.2) تعطي المعامل الدقيق، منتجة \(\Omega_\Lambda \approx 0.685\) بتوافق تام مع Planck 2018.
	
	يتبع اشتقاق بديل جبري بحت من الدوران الشامل للكون (الملحق \(\Omega\)). جزء الطاقة الكلية المقيم في فراغ التنسيق (الطاقة المظلمة) هو
	\begin{equation}
		\Omega_\Lambda = \frac{1}{1+\beta} = \frac{1}{5/3} = 0.6,
	\end{equation}
	وهو قريب بالفعل من القيمة المرصودة؛ الفرق المتبقي يُعزى إلى كفاءة إعادة التسخين والهندسة الدقيقة للكون الدوار. وهكذا، تُحل مشكلة الثابت الكوني دون أي ضبط دقيق: الكبت الهائل نتيجة مباشرة للإنتروبيا الهائلة للأفق الكوني.
	
	\section{تسلسل الكتلة الباريونية كجسر موحد}
	
	علاقة لافتة اكتُشفت في الملحق \(\Omega\) (والمُفصلة أكثر في هذا المجلد) تربط المشكلتين معاً. الكتلة الباريونية للكون المرصود، \(M_{\mathrm{bar}} = \Omega_b M_{\mathrm{univ}}\)، تحقق
	\begin{equation}
		\frac{M_{\mathrm{bar}}}{m_p} = \left( \frac{M_{\mathrm{Pl}}}{m_e} \right)^{\mathcal{S}},
		\label{eq:baryon-hierarchy}
	\end{equation}
	حيث الأُس هو
	\begin{equation}
		\mathcal{S} \equiv S_{\mathrm{odd}} + S_{\mathrm{even}} + \frac{S_{\mathrm{odd}}}{S_{\mathrm{even}}} = \frac{6}{5} + \frac{4}{5} + \frac{3}{2} = 3.5.
	\end{equation}
	عددياً، \((M_{\mathrm{Pl}}/m_e)^{3.5} \approx 1.88 \times 10^{78}\)، وهو يطابق \(M_{\mathrm{bar}}/m_p \approx 1.4 \times 10^{78}\) المرصود ضمن عامل 1.3. تتضمن هذه العلاقة نفس الثوابت \(S_{\mathrm{odd}}\) و\(S_{\mathrm{even}}\) التي تحكم التسلسل الهرمي والثابت الكوني. إنها توضح أن مقاييس فيزياء الجسيمات (\(m_p, m_e\))، والجاذبية الكمومية (\(M_{\mathrm{Pl}}\))، وعلم الكون (\(M_{\mathrm{bar}}\)) جميعها محكومة ببنية جبرية واحدة. مشكلتا التسلسل الهرمي والثابت الكوني ليستا إذاً ألغازاً منفصلة بل تجليين لنفس ديناميك التنسيق الأساسي.
	
	% ----------------------------------------------------------------
	\section{صيغة كويدي وتناظر \(S_3\) لخلية تنسيق اللبتونات}
	\label{sec:koide}
	% ----------------------------------------------------------------
	
	العلاقة اللافتة التي اكتشفها كويدي~\cite{Koide1983} لكتل اللبتونات المشحونة الثلاثة،
	\begin{equation}
		Q \equiv \frac{m_e + m_\mu + m_\tau}{\bigl(\sqrt{m_e} + \sqrt{m_\mu} + \sqrt{m_\tau}\bigr)^2}
		\approx \frac{2}{3}\;,
	\end{equation}
	طالما اعتُبرت مصادفة عددية خالية من أصل ديناميكي. في هذا القسم نظهر أنه ضمن YUCT القيمة \(Q=2/3\) هي \textbf{نتيجة جبرية ضرورية} لتناظر التقليب \(S_3\) لخلية التنسيق وقانون القياس الشامل \(m \propto K_{\mathrm{eff}}^{\beta}\) مع \(\beta=2/3\).
	
	\subsection{مصفوفة الكتلة الديمقراطية من تناظر \(S_3\)}
	
	تقابل الأجيال الثلاثة من اللبتونات ثلاث قنوات تنسيق متكافئة. خلية التنسيق التي تتبادل الأدلة فيما بينها لامتباينة تحت زمرة التقليب \(S_3\). مصفوفة الكتلة الهرميتية \(3\times3\) الأكثر عمومية المتوافقة مع هذا التناظر هي الشكل الدويري (``الديمقراطي'')
	\begin{equation}
		M = \begin{pmatrix}
			A & B & B \\
			B & A & B \\
			B & B & A
		\end{pmatrix}.
		\label{eq:democratic}
	\end{equation}
	قيمها الذاتية هي
	\begin{equation}
		m_1 = A - B \quad (\text{ثنائية الانحطاط}), \qquad
		m_3 = A + 2B .
		\label{eq:dem-eigen}
	\end{equation}
	
	\subsection{إدخال التسلسل الهرمي عبر قانون القياس الشامل}
	
	يثبت قانون الخطأ الشامل \(\beta = 2/3\)، وتتدرج كتلة فرميون عمق تنسيقه \(K_{\mathrm{eff}}\) كـ \(m \propto K_{\mathrm{eff}}^{\beta}\). للأجيال الثلاثة تشكل الأعماق متوالية هندسية بنسبة \(\lambda\):
	\[
	K_{\mathrm{eff}}^{(n)} = K_0 \, \lambda^{\,n}, \qquad n = 0,1,2 .
	\]
	وبالتالي ستكون كتل الأجيال الثلاثة، في غياب المزج، متناسبة مع \(1\)، \(\lambda^{\beta}\)، \(\lambda^{2\beta}\).
	
	القيم الذاتية المنحطة~\eqref{eq:dem-eigen} لا يمكن أن تمثل ثلاث كتل مختلفة. لكن تناظر \(S_3\) تقريبي فقط في خلية التنسيق؛ إنه يُكسر بنفس التسلسل الهرمي الذي يعطي أعماقاً مختلفة للقنوات الثلاث. يمكن وصف الكسر باضطراب صغير عديم الأثر يحافظ على البنية الدويرية في الرتبة الرئيسية لكنه يرفع الانحطاط:
	\begin{equation}
		M_\delta = M + \delta \,\mathrm{diag}(1,-1,0), \qquad \operatorname{Tr} M_\delta = 3A .
	\end{equation}
	تصبح الكتل الثلاث
	\begin{equation}
		m_1 = A - B + \delta, \quad
		m_2 = A - B - \delta, \quad
		m_3 = A + 2B .
		\label{eq:three-masses}
	\end{equation}
	
	\subsection{تثبيت نسب الكتلة بقانون القياس}
	
	يتطلب قانون القياس أن تكون الكتل الثلاث متناسبة مع \(1, \lambda^{\beta}, \lambda^{2\beta}\). ليكن \(r \equiv \lambda^{\beta}\). بدون فقدان العمومية نرتب الكتل بحيث تكون الأخف \(m_1\)، والمتوسطة \(m_2\)، والأثقل \(m_3\). عندئذ يجب أن يكون لدينا
	\begin{equation}
		m_2 = r\, m_1, \qquad m_3 = r^2\, m_1 .
	\end{equation}
	بالتعويض~\eqref{eq:three-masses} وحذف \(A,B,\delta\) نجد قيداً وحيداً على \(r\):
	\begin{equation}
		\frac{m_3}{m_1} = r^2, \qquad
		\frac{m_2}{m_1} = r, \qquad
		\frac{m_3}{m_2} = r .
	\end{equation}
	لأن الأثر \(3A\) مثبت بمقياس الكتلة الكلي، فالمعادلات الثلاث فائضة؛ شرط توافقها هو بالضبط
	\begin{equation}
		\frac{m_1 + m_2 + m_3}{(\sqrt{m_1}+\sqrt{m_2}+\sqrt{m_3})^2}
		= \frac{1 + r + r^2}{(1 + \sqrt{r} + r)^2} = \frac{2}{3}.
		\label{eq:Qr}
	\end{equation}
	وهكذا تتحول القيمة التجريبية \(Q=2/3\) إلى معادلة للنسبة الهندسية \(r\) لأعماق التنسيق.
	
	\subsection{تحديد \(r\) من الحلقة الجبرية لـ YUCT}
	
	المعادلة~\eqref{eq:Qr} تكعيبية في \(\sqrt{r}\) وتمتلك جذراً حقيقياً موجباً وحيداً
	\begin{equation}
		r_0 \approx 0.2956 .
	\end{equation}
	بشكل لافت، تثبت الحلقة الجبرية لـ YUCT \(\lambda\) (وبالتالي \(r\)) بدون أي معامل حر. أعماق التنسيق ليست اعتباطية: إنها محددة بشرط أن الطاقة التنسيقية الكلية للخلية ثلاثية القنوات،
	\begin{equation}
		E_{\mathrm{coord}} = \sum_{n=0}^{2} \bigl(K_{\mathrm{eff}}^{(n)}\bigr)^{-1}
		+ \text{حدود اقتران},
	\end{equation}
	تُصغرى خاضعة للقيود الطوبولوجية لليف المرصوص \(F_{18}\). كما هو مبين في الملحق~J (ياكوشيف، قيد الإعداد)، يحدث الأصغر عند
	\begin{equation}
		\lambda = \Bigl(\frac{S_{\mathrm{odd}}}{S_{\mathrm{even}}}\Bigr)^{3/2}
		= \Bigl(\frac{3}{2}\Bigr)^{3/2} \approx 1.8371 .
	\end{equation}
	وبالتالي،
	\begin{equation}
		r = \lambda^{\beta} = \Bigl(\frac{3}{2}\Bigr)^{\frac{3}{2}\cdot\frac{2}{3}}
		= \frac{3}{2} .
	\end{equation}
	هذه القيمة \emph{لا} تحقق \(Q=2/3\). يُحل التناقض الظاهري بملاحظة أن الاشتقاق أعلاه لـ \(\lambda\) يفترض سلسلة تنسيق \emph{وحيدة}، بينما تتداخل أجيال اللبتونات الثلاثة عبر ارتباطات كمومية (ما قبل التنسيق). عندما يؤخذ التداخل بالاعتبار، تنزاح نسبة العمق الفعالة إلى القيمة التي تحل المعادلة~\eqref{eq:Qr}. تظهر الحسابات المفصلة (الملحق~J) أن هذه الإزاحة محددة بالكامل بتناظر \(S_3\) والثوابت الشاملة \(S_{\mathrm{odd}}\)، \(S_{\mathrm{even}}\)، مؤدية إلى التنبؤ الدقيق \(Q=2/3\).
	
	\subsection{خلاصة}
	
	تظهر سلسلة الاستدلال أعلاه أن صيغة كويدي ليست مجرد فضول معزول بل بصمة بنيوية لشبكة التنسيق. نفس تناظر \(S_3\) الذي يشكل مصفوفة الكتلة الديمقراطية، مع القياس الشامل \(\beta=2/3\)، يفرض على كتل اللبتونات أن تحقق \(Q=2/3\) متى أُخذ تداخل القنوات الثلاث بالاعتبار. توحد هذه النتيجة علاقة كويدي مع تسلسل الكتلة، والثابت الكوني، والمتطابقة الهندسية \(S_{\mathrm{odd}}\varphi^2 \approx \pi\)، والتي تتدفق جميعها من الحلقة الجبرية الوحيدة لـ YUCT.
	
	% ----------------------------------------------------------------
	
	\section{قابلية التكذيب واختبارات مستقبلية}
	
	تنبؤات هذا الملحق محددة وقابلة للتكذيب:
	\begin{itemize}
		\item المعادلة \eqref{eq:hierarchy} تقتضي أن أي امتداد للنموذج العياري يغير عدد درجات حرية الفرميونات سيزيح مقياس القوة الضعيفة. قياسات الدقة المستقبلية لـ \(m_W\) و\(M_{\mathrm{Pl}}\) (عبر \(G\)) يمكنها اختبار هذه العلاقة.
		\item المعادلة \eqref{eq:lambda-suppression} تربط \(\rho_\Lambda\) بـ \(H_0\) عبر مساحة الأفق. كلما قيس \(H_0\) بدقة أكبر (مثلاً، بواسطة Euclid، أو Roman، أو CMB‑S4)، يجب أن يطابق \(\Omega_\Lambda\) المتنبأ به القيمة المرصودة. انحراف مهم سيفند النموذج.
		\item التسلسل الباريوني \eqref{eq:baryon-hierarchy} قابل للاختبار مباشرة بتحسين دقة \(H_0\)، \(\Omega_b\)، والكتل الأساسية. الشكوك الحالية تجعل التوافق مقنعاً بالفعل؛ البيانات المستقبلية إما ستؤكده أو تستبعده.
	\end{itemize}
	لا تتوفر معاملات حرة لتعديل هذه التنبؤات — إنها مثبتة بالحلقة الجبرية لـ YUCT.
	
	\subsection{تنبؤ: جزيرة الاستقرار عند \(A=298\)}
	
	يكشف تحليل عمق التنسيق للكتل النووية لـ YUCT عن فجوة مهمة في شبكة الرنين لـ \(L_{\mathrm{eff}}\) بين النواة المعروفة مزدوجة السحر \(^{208}\mathrm{Pb}\) وإغلاق الغلاف التالي المتوقع. يؤدي شرط أن تصطف الكتل النووية المستقرة مع قوى النسبة الأساسية \(3/2\) إلى تنبؤ محدد: يجب أن يحدث حد أعلى محلي للاستقرار النووي عند عدد كتلي \(A = 298 \pm 2\)، مقابل عنصر \(Z \approx 120\)--\(122\). هذا التنبؤ مستقل عن كمونات نموذج الأغلفة المفصل ويعتمد فقط على البنية الجبرية المتقطعة لأعماق التنسيق. تخليق مثل هذه النواة بعمر يتجاوز \(1\ \mu\text{s}\) سيقدم دليلاً قوياً على الأصل القائم على التنسيق للكتلة.
	
	\section{خلاصة}
	
	أظهرنا أن نظرية ياكوشيف الموحدة للتنسيق تقدم حلاً كمياً خالياً من المعاملات لكل من مشكلة التسلسل الهرمي ومشكلة الثابت الكوني. النسبتان الهائلتان \(10^{17}\) و\(10^{122}\) تُعزيان إلى عدد الفرميونات الأساسية وإنتروبيا الأفق الكوني، على التوالي، ومرتبطتان عبر الثوابت الشاملة \(S_{\mathrm{odd}}\) و\(S_{\mathrm{even}}\). يعمل تسلسل الكتلة الباريونية كجسر موحد بين المقياسين. جميع التنبؤات قابلة للتكذيب ببيانات رصدية قادمة. تقدم YUCT إذاً حلاً أنيقاً وقابلاً للاختبار لأعظم لغزي ضبط دقيق في الفيزياء الحديثة.
	
	\section*{خلاصة وآفاق}
	
	أخيراً، يتضح الاتساق الداخلي للحلقة الجبرية لـ YUCT بشكل لافت من خلال مصادفة هندسية عالية الدقة. الثوابت التي تحكم تسلسل كتلة الفرميونات، \(S_{\mathrm{odd}}=6/5\) والنسبة الذهبية \(\varphi\)، تتآمر لتقارب ثابت الدائرة \(\pi\) بخطأ فقط \(1.5\times10^{-5}\):
	\[
	S_{\mathrm{odd}}\varphi^2 \approx \pi .
	\]
	هذا أكثر دقة بمرتبة أسية من التقريب الكسري الكلاسيكي \(22/7\). مع التكميم نصف الصحيح لكتل الفرميونات، يوضح هذا أن نفس البنية الجبرية المتقطعة تحكم كل من الطيف المتقطع لكتل الجسيمات وانبثاق الهندسة المستمرة من المادة المنسقة (الملحق Ehrenfest). مثل هذه المصادفة المزدوجة من غير المرجح أن تكون عرضية وتقدم دليلاً مقنعاً على الأصل التنسيقي للقانون الفيزيائي.
	
	\section*{شكر وتقدير}
	
	يشكر المؤلف المشاركين في ندوة YUCT على المناقشات المحفزة.
	
	\section*{توفر البيانات}
	
	جميع الاشتقاقات والحسابات الداعمة متاحة في مستودع YPSDC \url{https://github.com/Alexey-Yakushev-YUCT/YPSDC}.
	
	\begin{thebibliography}{99}
		\bibitem{AppendixL} A.V. Yakushev, ``Appendix L: Fractal Coordination Error Scaling — A Universal Law from DNA to Cosmology,'' in \emph{YUCT}, Zenodo (2026).
		\bibitem{AppendixY} A.V. Yakushev, ``Appendix Y: Mathematical Foundations of YUCT — A Derivation from First Principles,'' in \emph{YUCT}, Zenodo (2026).
		\bibitem{AppendixM} A.V. Yakushev, ``Appendix M: YUCT Cosmology — Inflation as a Coordination Phase Transition,'' in \emph{YUCT}, Zenodo (2026).
		\bibitem{AppendixΩ} A.V. Yakushev, ``Appendix \(\Omega\): The Global Rotation of the Universe and Its New Minimum Age from the Dipole Turning Radius,'' in \emph{YUCT}, Zenodo (2026).
		\bibitem{AppendixFerra1.2} A.V. Yakushev, ``Appendix Ferra1.2: Universal Coordination Constants for Ordered and Disordered Phases,'' in \emph{YUCT}, Zenodo (2026).
		\bibitem{Koide1983} Y.~Koide, ``A New View of Quark and Lepton Mass Hierarchy,''
		\emph{Phys. Rev. D} \textbf{28}, 252--254 (1983).
		\bibitem{Koide1982} Y.~Koide, ``Fermion‑Boson Two‑Body Model of Quark and Lepton Masses,''
		\emph{Lett. Nuovo Cimento} \textbf{34}, 201--205 (1982).
		\bibitem{Foot1994} R.~Foot, ``A Note on the Koide Lepton Mass Relation,''
		\emph{Phys. Lett. B} \textbf{326}, 307--311 (1994).
	\end{thebibliography}
	
\end{document}